% Adaptación del preámbulo de Paper III v1.5 español; no plantilla standalone de Pandoc.
\documentclass[english,11pt,a4paper]{article}
\usepackage[margin=1in]{geometry}
\usepackage{fontspec}
\setmainfont{lmroman10-regular.otf}[BoldFont=lmroman10-bold.otf,ItalicFont=lmroman10-italic.otf,BoldItalicFont=lmroman10-bolditalic.otf]
\setsansfont{lmsans10-regular.otf}[BoldFont=lmsans10-bold.otf]
\setmonofont{lmmono10-regular.otf}[Scale=MatchLowercase]
\newlength{\LeanCharWidth}
\usepackage{amsmath,amssymb}
\usepackage[english]{babel}
\usepackage{microtype,parskip}
\usepackage{graphicx,float,caption}
\usepackage{longtable,booktabs,array,calc}
\usepackage{xcolor}
\usepackage{xurl}
\usepackage{hyperref,bookmark}
\usepackage{etoolbox,needspace}
\setcounter{secnumdepth}{-\maxdimen}
\setlength{\emergencystretch}{3em}
\providecommand{\tightlist}{\setlength{\itemsep}{0pt}\setlength{\parskip}{0pt}}
\newcounter{none}
\AtBeginEnvironment{longtable}{\Needspace{8\baselineskip}\small}
\pretocmd{\section}{\Needspace{5\baselineskip}}{}{}
\pretocmd{\subsection}{\Needspace{4\baselineskip}}{}{}
\pretocmd{\subsubsection}{\Needspace{4\baselineskip}}{}{}
\captionsetup{font=small,labelfont=bf}
\urlstyle{same}
\hypersetup{colorlinks=true,linkcolor=blue,urlcolor=blue,citecolor=blue,pdflang={en-US}}
\title{Formal companion to the Erdős 81 transfer note}
\author{Juan Pablo Traverso Gianini}
\date{}\begin{document}\raggedbottom\maketitle
Formal assembly of the IV--V application\\
Preprint 1.0, October 18, 2026

This companion describes the final Lean source package for Series
research note I, Preprint 1.0. The note reorganizes part of the proof in
Paper IV through Paper V's transfer interface. It claims no new solution
of Erdős Problem 81 and no mathematically independent proof. The
statements and proof terms retain their established scope. The current
validation is bound to the source manifest and final logs identified
below.

\subsection{1. What the formalization
proves}\label{what-the-formalization-proves}

The public conclusions produce actual clique partitions of the original
graph. For all sufficiently large chordal graphs of order \(n\), the
partition has at most

\[
M(n)=\left\lfloor\frac{n(n+1)}6\right\rfloor
\]

blocks, each with at most four vertices. For all orders, including zero,
an absolute additive constant suffices. There is no assertion of the
sharp bound at every order, no numerical upper bound on the eventual
threshold, and no new extremal value claimed by this application.

The source root is \nolinkurl{Erdos81Note.lean}. It imports
\nolinkurl{Erdos81Note.Assembly}, which imports the physical transfer
adapter and the required results of IV directly. The assembly invokes
the actual realizable-class transfer theorem from V. It does not invoke
IV's final all-chordal bound. The project remains mathematically
dependent on the structural and rounding results of IV and on the
transfer theorem of V; the separate project records a different assembly
of those proved inputs. Paper VI is not an input.

The principal note keeps one paragraph on formalization. This companion
gives the model boundary, the six checked public interfaces, the
selected dependency chain, the complete local-project axiom audit, and
reproduction information.

\subsection{2. Graphs, resources and integer
arithmetic}\label{graphs-resources-and-integer-arithmetic}

The physical graph model is \texttt{SimpleGraph\ V} with a finite vertex
type, decidable equality and decidable adjacency. The first two public
interfaces work on an arbitrary finite type. The final graph theorems
use \texttt{Fin\ n}. Chordality is the predicate
\texttt{PaperIV.IsChordal\ G}, which is printed as \nolinkurl{G.IsChordal}
in the endpoint gate.

A packing item is a finite set of vertices that is a clique and has
cardinality three or four. Its resources are its unordered pairs of
distinct vertices, hence actual graph edges. A packing consists of such
items with disjoint edge supports. Its natural-valued gain is

\[
g(P)=\sum_{Q\in P}\left(\binom{|Q|}{2}-1\right).
\]

The adapter \nolinkurl{physicalCS} makes two distinct items conflict
precisely when their edge supports overlap. The maps \nolinkurl{toPacking}
and \nolinkurl{liftPacking} prove the correspondence between feasible
outputs of V and physical packings in IV; their gain lemmas preserve the
weighted sum. This step rules out a certificate whose alleged
realization is merely an auxiliary object without graph edges.

A \texttt{CliquePartition\ G} contains a finite family of cliques, each
of order at least two, with pairwise disjoint edge supports whose union
is exactly \nolinkurl{G.edgeFinset}. \nolinkurl{Q.size} counts blocks, and
\texttt{Q.OrderAtMost\ 4} bounds their orders. Completion adds one
two-vertex block for each uncovered edge and proves

\[
|Q|+g(P)=e(G).
\]

Thus the subtraction used in the paid-packing interface is natural
subtraction with a proved bound \(g(P)\le e(G)\); it is not an
unexamined truncation. Likewise, \texttt{targetSize\ n} is natural
division \(n(n+1)/6\), exactly \(M(n)\).

The rational predicate \texttt{CertifiedFractionalOptimum\ G\ w} means
that there exist a feasible fractional mixed packing and a feasible dual
cover, both of value \(w\). It supplies matching primal and dual
witnesses, rather than an unsupported assertion that a maximum exists.
The assembly obtains this certificate from IV's proved existence theorem
before splitting into the far and near cases.

\subsection{3. The six checked public
interfaces}\label{the-six-checked-public-interfaces}

All names in this section belong to the namespace \nolinkurl{Erdos81Note}.
The endpoint gate checks their full elaborated types and prints their
axiom dependencies. The final output is retained as
\nolinkurl{evidence/ENDPOINTS.log}; the following descriptions retain every
substantive hypothesis and conclusion.

\subsubsection{\texorpdfstring{3.1.
\nolinkurl{partition_of_certificate}}{3.1. partition\_of\_certificate}}\label{partition_of_certificate}

Let \(G\) be any finite simple graph, \(C\) an auxiliary conflict
system, and let the weight function \(w\), multiplicity function \(X\),
class count \(q\), and target \(T\) be natural-valued. Assume \(1\le q\)
and a certificate

\[
\texttt{RealCert}\bigl(C,\operatorname{physicalCS}(G),
w,\operatorname{physicalWeight}(G),X,q\bigr).
\]

The certificate includes deleted occurrences, an assignment of the
remaining occurrences to classes, and a feasible physical realizer for
each class with an explicitly bounded natural-valued loss. Its cost is
the deleted weight plus the sum of realization losses. If

\[
q\,e(G)+\operatorname{cost}(\mathrm{cert})
\le qT+\operatorname{wMultSize}(C,w,X),
\]

then there exists a \texttt{CliquePartition\ G} of order at most four
and size at most \(T\). No chordality hypothesis is needed. The proof
calls \nolinkurl{PaperV.realizable_class_transfer}, converts the selected
realization into a physical packing, and completes it.

\subsubsection{\texorpdfstring{3.2.
\nolinkurl{partition_of_paid_packing}}{3.2. partition\_of\_paid\_packing}}\label{partition_of_paid_packing}

For any finite graph \(G\), natural target \(T\), and physical mixed
packing \(P\), assume

\[
e(G)-g(P)\le T.
\]

Then there exists a clique partition of \(G\), of order at most four and
size at most \(T\). Internally, the proof chooses one auxiliary
occurrence with reference weight \(e(G)\), one class, and realization
loss \(e(G)-g(P)\), and applies the preceding interface. The
construction records a paid physical packing in the transfer language;
it does not claim an additional saving from averaging.

\subsubsection{\texorpdfstring{3.3.
\nolinkurl{far_certificate}}{3.3. far\_certificate}}\label{far_certificate}

For every rational \(\eta>0\), there exists a natural threshold \(N\)
such that, for every \(n\ge N\), every graph \(G\) on \texttt{Fin\ n},
and every rational certified optimum \(w\), the inequality

\[
e(G)-w<\frac{n^2}{6}-\eta n^2
\]

implies a clique partition of order at most four and size at most
\(M(n)\). This interface does \textbf{not} require chordality. It uses
IV's uniform transfer with parameter \(\eta/2\); the far-regime margin
pays the rounding loss. Its threshold is existential, not the explicit
near-branch threshold below.

\subsubsection{\texorpdfstring{3.4.
\nolinkurl{near_certificate}}{3.4. near\_certificate}}\label{near_certificate}

Let \(n\ge4\cdot10^{12}\), let \(G\) on \texttt{Fin\ n} be chordal, and
let \(w\in\mathbb Q\) have a certified fractional optimum witness.
Assume the negation of the strict far inequality:

\[
\neg\left(e(G)-w<\frac{n^2}{6}-\eta n^2\right),
\qquad \eta=10^{-16}.
\]

Then there exists a clique partition of order at most four and size at
most \(M(n)\). The proof obtains an actual near-structure witness
through IV's descent localization, uses its physical packing and
retained accounts, proves that its completion is paid, and invokes
\nolinkurl{partition_of_paid_packing}. Neither a target partition bound
nor a pre-existing root witness is assumed in this interface.

\subsubsection{\texorpdfstring{3.5.
\nolinkurl{erdos81_eventual}}{3.5. erdos81\_eventual}}\label{erdos81_eventual}

There exists \(N\in\mathbb N\) such that, for every \(n\ge N\) and every
chordal graph \(G\) on \texttt{Fin\ n}, there exists a clique partition
of order at most four and size at most \(M(n)\). The certified optimum
and the case split are obtained internally. The chosen threshold is the
maximum of the far-branch threshold at \(\eta=10^{-16}\) and
\(4\cdot10^{12}\); the latter alone is not asserted to suffice for the
global theorem.

\subsubsection{\texorpdfstring{3.6.
\nolinkurl{erdos81_all_orders}}{3.6. erdos81\_all\_orders}}\label{erdos81_all_orders}

There exists \(b\in\mathbb N\) such that, for every \(n\in\mathbb N\)
and every chordal graph \(G\) on \texttt{Fin\ n}, there exists a clique
partition of order at most four and size at most \(M(n)+b\). The proof
takes \(b=N^2\) from the previous theorem. Below the threshold it uses
individual edges; above it, the eventual theorem. No positive-order
assumption excludes the empty graph.

\subsection{4. Which results are used, and which are
excluded}\label{which-results-are-used-and-which-are-excluded}

The far branch calls \nolinkurl{PaperIV.RC01FarAssembly.uniformTransferAt}.
The near branch calls
\nolinkurl{PaperIV.NearH1Localization.exists_nearStructureWitness_of_nearRegime}
and reaches the descent-localization declaration
\nolinkurl{PaperIV.LocalStability.localize_original_of_windowed_accounts_descent}.
Optimum existence is supplied by
\nolinkurl{PaperIV.CertifiedOptimumExistence.exists_certifiedFractionalOptimum}.

The selected proof-chain audit begins with the certificate interface and
the two final bounds. It follows declaration types and proof values
through the local series and records 6,528 visited declarations. It
requires actual dependencies on V's transfer and IV's descent
localization. It rejects declarations originating in
\nolinkurl{PaperIV.Erdos81Unconditional}, the final global-bound module.
The recorded forbidden list and normalized first-entry hits are empty.

This is a dependency claim about the selected proof, not a claim that
all imported source files omit alternative arguments. Some upstream
files contain unused first-entry developments. Their presence in a
source package does not make those declarations premises of the final
application. The complete dependency listing is preserved as
\nolinkurl{DECLARATION_DEPENDENCIES.txt}.

The physical partition is constructed once in the original graph after
localization. The formalization does not transport partitions backwards
along vertex-copy steps and does not charge an extra loss per step. The
certificate theorem organizes the gain budget; the structural
localization and physical construction remain the work of IV.

\subsection{5. Final validation and trusted
boundary}\label{final-validation-and-trusted-boundary}

The canonical project uses \nolinkurl{leanprover/lean4:v4.35.0} and Mathlib
commit \nolinkurl{56a005855ea17e0b0958f73ec801ae140b0a1679}. The source
closure contains 324 local modules. The final validation consists of a
root build, the six-interface gate, and a full local-project audit. The
build uses complete, trace-checked cached artifacts for unchanged
dependencies and recompiles the changed modules; it is not described as
a fresh build of all upstream sources or of Mathlib.

The selected root audit traverses 6,528 declaration dependencies. It
checks actual use of V's transfer and IV's descent localization and
excludes IV's final global-bound module. The full audit covers 6,850
local declarations and visits 10,304 dependencies. These are distinct
scopes: the latter includes the private helpers and other declarations
in all imported local modules.

The full audit first verifies that every module in the source manifest
is present in the imported Lean environment; a missing module causes
rejection.

The permitted axiom union is \nolinkurl{propext}, \nolinkurl{Classical.choice}
and \nolinkurl{Quot.sound}. The gates reject additional mathematical axioms
and \nolinkurl{sorryAx}. The successful command exits and logs are recorded
in \nolinkurl{evidence/BUILD_CHECKS.json};
\nolinkurl{FINAL_FORMAL_VALIDATION.json} binds those records to the exact
source closure and the checks of the public entry point.

The author of the note's formalization is Juan Pablo Traverso Gianini.
Its own modules and gates identify the author in their headers.
\nolinkurl{AUTHORSHIP_NOTICE.md} credits the inherited series source cuts
without altering their existing mathematical and assistance
attributions. The Paper V dependencies use the author-credited source
version recorded in the manifest. The note's own modules and audit gates
are distributed under Apache License 2.0, with the complete license in
\nolinkurl{lean435/LICENSE} and source attribution in
\nolinkurl{lean435/NOTICE}. Inherited license and authorship notices remain
applicable.

\nolinkurl{lean435/OWNERSHIP_AND_LICENSE.json} enumerates, with hashes,
all author-owned Lean sources, including the reused series cuts,
released by the author under Apache 2.0. The new full headers cover this
note's specific modules and audit entry points; inherited headers and
mathematical attributions are preserved. External dependencies retain
their own licenses.

\Needspace{12\baselineskip}
\subsection{6. Reproduction and immutable
package}\label{reproduction-and-immutable-package}

\Needspace{8\baselineskip}
Extract \nolinkurl{NOTE_I_LEAN435_FINAL_SOURCE_v1.0.zip}, enter
\nolinkurl{lean435}, and run:

\begin{verbatim}
lake exe cache get
lake build Erdos81Note
lake env lean -j 1 -s 32768 Endpoint435Gate.lean
lake env lean -j 1 -s 32768 FreezeFullAudit.lean
\end{verbatim}

Preserve the supplied \nolinkurl{lake-manifest.json}; the cache command
retrieves artifacts for the pinned dependencies without requesting a new
dependency selection.

The root build runs the selected dependency audit in
\nolinkurl{Erdos81Note.lean}. The endpoint gate prints the six public types
and their axiom dependencies. The full audit checks every declaration in
the imported local source closure. The archive includes all required
local source files and pinned configuration; it excludes \nolinkurl{.lake}
caches and development junctions. The package does not require a
checkout of Papers IV or V.

\Needspace{8\baselineskip}
The source archive SHA-256 is:

\begin{verbatim}
eae69a0316113f1701246b54243a1b51696b85ea656309319b8738932f90f604
\end{verbatim}

\Needspace{8\baselineskip}
\nolinkurl{SOURCE_MANIFEST.json} gives the hash of every packaged source
and configuration file. Its own SHA-256 is:

\begin{verbatim}
5f7deb401e26fcce100d576e4ed5c6975a46d14265144a1f267439b67f24613f
\end{verbatim}

\nolinkurl{FINAL_FORMAL_VALIDATION.json} records the build, export and
axiom results for this package. Exact-arithmetic regression scripts are
supplementary checks, not premises of the universal theorems.

\subsection{Assistance and scope}\label{assistance-and-scope}

The principal note retains Paper IV's attribution of Claude, by
Anthropic, ChatGPT/Codex, by OpenAI, Aristotle, by Harmonic, Certo and
Jacobian for the inherited work. Aristotle's contribution to that work
included exploration, counterexample search, proof development,
formalization and review, and was not limited to translating finished
proofs. For this note, ChatGPT/Codex assisted with the adapters,
verification and presentation of the final package. Computations were
reviewed within their stated scope; a finite search does not replace a
universal proof. Kernel verification does not establish historical
priority or replace external review.

\Needspace{3\baselineskip}The author remains responsible for the arguments, citations, code, and
presentation. No AI system is listed as an author.

\end{document}
